\section{Repository und Reproduzierbarkeit}\label{app:repro}

Der gesamte Code liegt unter \url{https://github.com/yarx/JassCardEye}.
Die wichtigsten Verzeichnisse:

\begin{description}
  \item[Dataset] \path{src/tools/dataset/}: .NET-Lösung mit Domänenmodell, Render-Engine, CLI, Viewer und Checks.
  \item[Daten] \path{data/}: 72 Kartenscans, 25 reale Negativfotos und das reale Validierungsset mit Bildern und Quell-Labels.
  \item[\code{src/training/}] Trainings-, Export- und Fetch-Skripte (Python, Ultralytics) sowie die Pipeline-Dokumentation.
  \item[Pipeline] \path{src/scripts/} und \path{docker/} enthalten Laufskript und Container. Die Workflows liegen unter \path{.github/workflows/}. Der Release-Workflow baut beide Apps mit demselben Modell.
  \item[Apps] \path{src/app/ios/} und \path{src/app/android/}: SwiftUI und Kotlin.
  \item[\code{context/}] Projektwissen: Vision, Architektur inklusive der Tabelle aller Werte zwischen Kamera und gezählter Karte.
  \item[\code{doc/}] Diese Arbeit und das Arbeitsjournal (LaTeX).
\end{description}

Datensatz und Training lassen sich mit drei Befehlen reproduzieren. Derselbe
Seed erzeugt auf jeder Maschine bit-identische Bilder. Die handgelabelten
Validierungsdaten sind versioniert, synthetische Daten und Trainingsartefakte
werden bewusst erzeugt statt eingecheckt (Abschnitt~\ref{sec:repro}).

\begin{verbatim}
dotnet run -c Release --project src/tools/dataset/Cli -- \
    generate --count 200000 --seed 1
dotnet run -c Release --project src/tools/dataset/Cli -- \
    rebuild --dataset data/real/val
src/training/.venv/bin/python src/training/train.py --variant c
\end{verbatim}

Ein vollständiger Cloud-Lauf wird mit \code{src/scripts/run\_training.sh
--variant all} oder über den Workflow \code{train} gestartet. Jeder Lauf
liegt unter \code{training-runs/<Datum-Commit>/} in Azure Blob Storage mit
einem Manifest, das Commit, Seed, Bildanzahl, Hardware und Metriken nennt.
Das veröffentlichte Repository enthält den Endstand des Projekts ohne die
Entwicklungsgeschichte. Die Manifeste der Trainingsläufe nennen deshalb
Commits des privaten Entwicklungsrepositorys. Generator, Kartenscans,
Negativfotos und Trainingsskript der Datenmengenreihe aus
Kapitel~\ref{sec:experimente} entsprechen dem veröffentlichten Stand.
\code{src/training/fetch\_run.py
--list} zeigt alle Läufe, \code{--export} erzeugt daraus die Modelle für
beide Apps samt \code{models.json}.
